\titolosezione{Come usare questo canzoniere}
\begingroup\small
\begin{multicols}{2}
\textbf{Struttura.} I canti sono numerati e raggruppati in sezioni, nell'ordine:
<<SEZIONI>>. Ogni canto ha la sua pagina: il numero grande accanto al titolo
è quello da cercare, sia qui sia nel sito. In fondo trovi l'indice alfabetico
e l'indice per sezione.

\medskip
\textbf{Accordi.} Sono scritti sopra la sillaba in cui cambiano. Il testo in
\textit{corsivo} è il ritornello.

\medskip
\textbf{Online.} Sull'ultima pagina c'è un QR code: porta al sito con l'audio di
ogni canto, gli spartiti in PDF e la modalità notte per il bivacco.

\columnbreak
\textbf{Branche.} A destra del titolo, le branche per cui il canto nasce o è
più adatto:

\smallskip
\bLC\ lupetti e coccinelle \quad \bEG\ esploratori e guide \quad \bRS\ rover e scolte

\medskip
\textbf{Intensità.} Quanta energia chiede il canto, da 1 a 5:

\smallskip
\fiamme{1}\ raccolto, quasi sussurrato\\
\fiamme{3}\ corale, a voce piena\\
\fiamme{5}\ esplosivo, con gesti e balli

\medskip
Un canto senza tag è per tutti.
\end{multicols}
\endgroup
